% !TeX program = xelatex
\documentclass[11pt,a4paper]{ctexart}
% 11/9 薄弱点定向加练：诊断卷 12 + 六模块 72 = 84 题
% 由 build_weakness_drill.py 在 sympy/mpmath 全 PASS 后自动生成，答案不可手改
\usepackage[a4paper,top=1.7cm,bottom=1.7cm,left=1.7cm,right=1.7cm]{geometry}
\usepackage{amsmath,amssymb}
\usepackage{booktabs,array,longtable}
\usepackage{enumitem}
\usepackage[table]{xcolor}
\usepackage{fancyhdr}
\usepackage{hyperref}
\hypersetup{hidelinks,colorlinks=true,linkcolor=blue!50!black}
\usepackage{titlesec}
\linespread{1.14}
\setlength{\parskip}{0.10em}
\setlist[itemize]{itemsep=1pt,topsep=2pt,parsep=0pt,leftmargin=1.6em}
\titleformat{\section}{\large\bfseries\color{blue!35!black}}{\thesection}{0.5em}{}
\titlespacing*{\section}{0pt}{1.0ex plus .2ex}{0.5ex}
\pagestyle{fancy}
\fancyhf{}
\fancyhead[L]{\small 薄弱点定向加练 · 11/9 · 84 题}
\fancyhead[R]{\small 陶理}
\fancyfoot[C]{\small\thepage}
\renewcommand{\headrulewidth}{0.4pt}
\newcommand{\eps}{\varepsilon}
\newcommand{\ee}{\mathrm{e}}
\newcommand{\dd}{\mathop{}\!\mathrm{d}}
\newcommand{\blank}[1]{\underline{\hspace{#1}}}
\newcommand{\qi}[3]{\par\noindent\textbf{#1}\ \ {\footnotesize[#2]}\quad #3\par\smallskip}
\begin{document}
\begin{center}
{\LARGE\bfseries 薄弱点定向加练}\\[0.3em]
{\large 2026-11-09（D-2）\quad 诊断卷 12 + 六模块 72 = \textbf{84 题}}\\[0.3em]
{\small 先诊断、后加权；全部答案经 sympy / mpmath 独立复核锁死}
\end{center}
\begin{center}
\fcolorbox{blue!40!black}{blue!4}{
\begin{minipage}{0.95\textwidth}
\vspace{0.25em}
\textbf{这一册为什么这样配题}\quad 诊断依据见 \texttt{Notes/薄弱点诊断\_2026-11-09.md}：
你自己连续两周写下「函数性态作图零训练 → 仍卡住」，另有 M 型放缩方向、换元条件漏写 7 处、泰勒限时卷未达 14/20 三条明确记录。
\textbf{A 级四条各 14 题（56 题，占 78\%）}，B 级两条各 8 题；难度按基础 / 进阶 / 压轴三档。
\vspace{0.25em}
\end{minipage}}
\end{center}
\section{先看这一条：全库 0 条做题回执}
\textbf{诊断不是正确率排序}：2024-25 旧卷得分与错题号无记录，11/8 的 8 组正确率表 / S4 独立做出题数 / 三类错误次数全部空白，11/8 日报把「D-3 到期零回执」标红为 P13。
所以下面的权重是「自述卡点 × 命题分值 × 反复复发」的合成，\textbf{不编造分数}。
\section{诊断卷 D1--D12（20 min 掐表，先做再对答案）}
六弱点各 2 题。只填下面 3 行式记录，不要写长复盘；填完把 0--1/2 的模块整组提前到今晚第一段。
\qi{D1}{作图七步·渐近线}{求 $y=\dfrac{x^{2}+2}{x-1}$ 的斜渐近线。}
\qi{D2}{作图七步·渐近线}{求 $y=\dfrac{x^{3}}{x^{2}-1}$ 的斜渐近线。}
\qi{D3}{第 11 题 M 型}{用定义证 $\lim\limits_{x\to0}\dfrac1{x^{2}}=+\infty$（按 $|f|\to\infty$ 写），给出 $\delta(M)$。}
\qi{D4}{第 11 题 $\eps$-$\delta$}{用 $\eps$-$\delta$ 定义证 $\lim\limits_{x\to2}(3x-1)=5$，给出 $\delta(\eps)$。}
\qi{D5}{条件书写·参数二阶}{设 $x=t^{2},\ y=t^{3}$，求 $\dfrac{\dd^{2}y}{\dd x^{2}}\big|_{t=1}$，并写出第一句必须声明的前提。}
\qi{D6}{条件书写·隐函数}{设 $x^{2}+xy+y^{2}=7$，求 $y\prime$ 在 $(1,2)$ 处的值，并写出隐函数求导的前提。}
\qi{D7}{泰勒·通项系数}{求 $f(x)=(1+x)^{1/4}$ 的 $f\prime\prime\prime(0)$。}
\qi{D8}{泰勒·展错阶}{求 $\lim\limits_{x\to0}\dfrac{\ln(1+x)-x+\tfrac12x^{2}}{x^{3}}$。}
\qi{D9}{极限计算·加减禁替换}{求 $\lim\limits_{x\to0}\dfrac{\tan x-\sin x}{x^{3}}$。}
\qi{D10}{间断·跳跃}{求 $f(x)=\dfrac{\ee^{1/x}-1}{\ee^{1/x}+1}$ 在 $x\to0^{-}$ 与 $x\to0^{+}$ 的极限，并判间断类型。}
\qi{D11}{新压轴·迭代}{设 $a_1=1,\ a_{n+1}=\dfrac{a_n}{1+2a_n}$，求 $\lim\limits_{n\to\infty} n a_n$。}
\qi{D12}{极坐标切线}{极坐标 $r=1+\cos\theta$，求 $\theta=\dfrac\pi2$ 处切线斜率。}

\begin{center}\small
\setlength{\tabcolsep}{3pt}
\begin{tabular}{@{}p{4.2cm}p{2.2cm}p{9.0cm}@{}}
\toprule
\textbf{模块} & \textbf{对 / 2} & \textbf{错的「第一步」（一句话）} \\
\midrule
作图七步 & \blank{1.6cm} & \blank{8.4cm} \\
定义 M 型 / 量词 & \blank{1.6cm} & \blank{8.4cm} \\
条件书写 & \blank{1.6cm} & \blank{8.4cm} \\
泰勒 & \blank{1.6cm} & \blank{8.4cm} \\
极限计算 / 间断 & \blank{1.6cm} & \blank{8.4cm} \\
新压轴 / 极坐标 & \blank{1.6cm} & \blank{8.4cm} \\
\bottomrule
\end{tabular}
\end{center}
\section{函数性态作图七步（14 题）}
\noindent{\small\color{blue!35!black} A 级第 1 名 · 连续两周「仍卡住」 · 约 11 分 / 16 套}\\[0.3em]
\noindent\textbf{方法卡}\ 七步：定义域 → 奇偶 → 截距 → $y\prime$ → 极值 → $y\prime\prime$ → 渐近线。两端渐近线都要查。\\[0.5em]
\qi{M1.1}{基础 · 斜渐近线}{求 $y=\dfrac{x^{2}+3}{x-1}$ 的斜渐近线。}
\qi{M1.2}{基础 · 斜渐近线}{求 $y=\dfrac{2x^{2}+3x+1}{x+1}$ 的斜渐近线。}
\qi{M1.3}{基础 · 斜+垂直}{求 $y=\dfrac{x^{3}}{x^{2}-4}$ 的斜渐近线与垂直渐近线。}
\qi{M1.4}{进阶 · 斜渐近线}{求 $y=\dfrac{x^{3}+2x}{x^{2}+1}$ 的斜渐近线。}
\qi{M1.5}{进阶 · 斜渐近线·指数}{求 $y=x\ee^{1/x}$ 的斜渐近线。}
\qi{M1.6}{进阶 · 两端渐近线}{求 $y=\sqrt{x^{2}+4}$ 在 $x\to+\infty$ 与 $x\to-\infty$ 的两条渐近线。}
\qi{M1.7}{压轴 · 两端渐近线}{求 $y=x+\arctan x$ 的两条斜渐近线。}
\qi{M1.8}{基础 · 驻点}{求 $y=x^{3}-3x^{2}+2$ 的驻点。}
\qi{M1.9}{基础 · 驻点}{求 $y=x^{2}\ee^{-x}$ 的驻点。}
\qi{M1.10}{基础 · 拐点}{求 $y=x^{3}-3x^{2}$ 的拐点横坐标。}
\qi{M1.11}{进阶 · 拐点·两处}{求 $y=x^{4}-2x^{3}$ 的拐点横坐标。}
\qi{M1.12}{压轴 · 极值}{求 $y=x^{2}\ln x\ (x>0)$ 的极小值。}
\qi{M1.13}{基础 · 垂直渐近线计数}{函数 $y=\dfrac1{(x^{2}-1)(x-2)}$ 有几条垂直渐近线？}
\qi{M1.14}{进阶 · 水平渐近线}{求 $y=\dfrac{2x^{2}+1}{x^{2}-3x+2}$ 的水平渐近线。}
\section{第 11 题：量词行、M 型与放缩方向（14 题）}
\noindent{\small\color{blue!35!black} A 级第 2 名 · 17/17 套必考 8 分 · 官方答案栽过 2 次}\\[0.3em]
\noindent\textbf{方法卡}\ $\eps$-$\delta$ 往小估（上界），M 型往大估（下界）；分母要上界 / 分子要下界的口诀分清用在哪一步。\\[0.5em]
\qi{M2.1}{基础 · $\eps$-$\delta$}{$\lim\limits_{x\to1}(3x-1)=2$ 的 $\delta(\eps)$。}
\qi{M2.2}{进阶 · 先关区间}{$\lim\limits_{x\to2}x^{2}=4$ 的 $\delta(\eps)$（先关区间）。}
\qi{M2.3}{进阶 · 先关区间}{$\lim\limits_{x\to3}(x^{2}-1)=8$ 的 $\delta(\eps)$。}
\qi{M2.4}{进阶 · 放缩方向}{$\lim\limits_{x\to1}\dfrac1x=1$ 的 $\delta(\eps)$。}
\qi{M2.5}{进阶 · 有界×无穷小}{$\lim\limits_{x\to0}x\sin\dfrac1x=0$ 的 $\delta(\eps)$。}
\qi{M2.6}{压轴 · 放缩·二阶余项}{$\lim\limits_{x\to0}\dfrac{\sin x}{x}=1$ 的一个可行 $\delta(\eps)$（用 $|\tfrac{\sin x}{x}-1|\le \tfrac{x^{2}}6$）。}
\qi{M2.7}{基础 · $\eps$-$X$}{$\lim\limits_{x\to+\infty}\dfrac1x=0$ 的 $X(\eps)$。}
\qi{M2.8}{基础 · $\eps$-$X$}{$\lim\limits_{x\to+\infty}\dfrac{x}{x+1}=1$ 的 $X(\eps)$。}
\qi{M2.9}{进阶 · $\eps$-$X$·先关区间}{$\lim\limits_{x\to+\infty}\dfrac{2x+1}{x-1}=2$ 的 $X(\eps)$。}
\qi{M2.10}{基础 · M 型}{$\lim\limits_{x\to0}\dfrac1{x^{2}}=+\infty$ 的 $\delta(M)$。}
\qi{M2.11}{进阶 · M 型}{$\lim\limits_{x\to0}\dfrac1{|x|^{3}}=+\infty$ 的 $\delta(M)$。}
\qi{M2.12}{进阶 · M 型·非零趋近点}{$\lim\limits_{x\to1}\dfrac1{|x-1|}=+\infty$ 的 $\delta(M)$。}
\qi{M2.13}{压轴 · M 型·对数}{$\lim\limits_{x\to0}\ln\dfrac1{x^{2}}=+\infty$ 的 $\delta(M)$。}
\qi{M2.14}{压轴 · 用定义证不存在}{用定义证 $\lim\limits_{x\to0}\cos\dfrac1x$ 不存在（给出两列）。}
\section{条件书写与六型求导（14 题）}
\noindent{\small\color{blue!35!black} A 级第 3 名 · 作业 12–14 漏写 7 处 · 导数块 16→20 分}\\[0.3em]
\noindent\textbf{方法卡}\ 每道题第一句写前提：$x\prime_t\ne0$、$F_y\ne0$、$f\prime\ne0$、真数为正、$|x|<1$。\\[0.5em]
\qi{M3.1}{基础 · 参数二阶}{$x=t^{2},y=t^{3}$，求 $\dfrac{\dd^{2}y}{\dd x^{2}}\big|_{t=1}$；前提：$x\prime_t\ne0$。}
\qi{M3.2}{基础 · 参数一阶}{$x=\ee^{t}\cos t,\ y=\ee^{t}\sin t$，求 $\dfrac{\dd y}{\dd x}\big|_{t=0}$；前提：$x\prime_t\ne0$。}
\qi{M3.3}{进阶 · 参数·化简}{$x=t-\ln(1+t),\ y=t^{2}$，求 $\dfrac{\dd y}{\dd x}\big|_{t=1}$；前提：$x\prime_t\ne0,\ t>-1$。}
\qi{M3.4}{基础 · 隐函数一阶}{$x^{2}+xy+y^{2}=7$ 在 $(1,2)$ 求 $y\prime$；前提：$F_y\ne0$。}
\qi{M3.5}{进阶 · 隐函数·先验点}{$x^{3}+y^{3}=3xy$ 在 $(\tfrac32,\tfrac32)$ 求 $y\prime$；前提：$F_y\ne0$。}
\qi{M3.6}{进阶 · 隐函数·指数}{$\ee^{y}+xy=\ee$ 在 $(0,1)$ 求 $y\prime$；前提：$F_y\ne0$。}
\qi{M3.7}{压轴 · 反函数二阶}{$f(t)=t^{3}+t$，$y=f^{-1}(x)$，求 $y\prime\prime(0)$；前提：$f\prime(0)\ne0$。}
\qi{M3.8}{进阶 · 对数求导}{$y=(1+x)^{\sin x}$，求 $y\prime(0)$；前提：$1+x>0$。}
\qi{M3.9}{基础 · 分段可导·两条件}{$f(x)=\begin{cases}ax+b,&x\le0\\ \ee^{2x},&x>0\end{cases}$ 在 $x=0$ 可导，求 $a,b$。}
\qi{M3.10}{进阶 · 分段可导}{$f(x)=\begin{cases}x^{2}+c,&x\le1\\ dx+2,&x>1\end{cases}$ 在 $x=1$ 可导，求 $c,d$。}
\qi{M3.11}{基础 · Leibniz 高阶导}{求 $f(x)=x^{2}\ee^{x}$ 的 $f^{(5)}(0)$。}
\qi{M3.12}{进阶 · 参数·分子为零}{$x=\cos t,\ y=\sin 2t$，求 $\dfrac{\dd y}{\dd x}\big|_{t=\pi/4}$；前提：$x\prime_t\ne0$。}
\qi{M3.13}{基础 · 微分·前提}{写出 $\dd y$：$y=\ln|x|$，并声明前提。}
\qi{M3.14}{进阶 · 复合求导·前提}{$y=\arcsin(x^{2})$，求 $y\prime$ 并声明定义域前提。}
\section{泰勒：展开阶、系数与符号（14 题）}
\noindent{\small\color{blue!35!black} A 级第 4 名 · 限时卷未达 14/20 · 2023 起稳定大题}\\[0.3em]
\noindent\textbf{方法卡}\ $f^{(k)}(0)=k![x^k]$；$(1+x)^\alpha$ 的通项系数不要再乘 $k!$；分母几阶就展到几阶。\\[0.5em]
\qi{M4.1}{基础 · 通项系数}{求 $f(x)=(1+x)^{1/4}$ 的 $f\prime\prime\prime(0)$。}
\qi{M4.2}{基础 · 通项系数}{求 $f(x)=(1+x)^{-1/2}$ 的 $f^{(4)}(0)$。}
\qi{M4.3}{基础 · 复合展开}{求 $f(x)=\ee^{x^{2}}$ 的 $f^{(6)}(0)$。}
\qi{M4.4}{进阶 · 复合展开·符号}{求 $f(x)=\sin(x^{2})$ 的 $f^{(6)}(0)$。}
\qi{M4.5}{基础 · 级数乘移位}{求 $f(x)=\dfrac{x^{2}}{1-x}$ 的 $f^{(5)}(0)$。}
\qi{M4.6}{基础 · 展错阶}{求 $\lim\limits_{x\to0}\dfrac{\cos x-1+\tfrac12x^{2}}{x^{4}}$。}
\qi{M4.7}{基础 · 展错阶}{求 $\lim\limits_{x\to0}\dfrac{\tan x-x}{x^{3}}$。}
\qi{M4.8}{基础 · 展错阶}{求 $\lim\limits_{x\to0}\dfrac{\arcsin x-x}{x^{3}}$。}
\qi{M4.9}{基础 · 展错阶}{求 $\lim\limits_{x\to0}\dfrac{\ee^{x}-1-x-\tfrac12x^{2}}{x^{3}}$。}
\qi{M4.10}{进阶 · 展错阶·符号}{求 $\lim\limits_{x\to0}\dfrac{\ln(1+x)-x+\tfrac12x^{2}-\tfrac13x^{3}}{x^{4}}$。}
\qi{M4.11}{进阶 · 复合·高阶}{求 $f(x)=\sin(x^{3})$ 的 $f^{(9)}(0)$。}
\qi{M4.12}{压轴 · 逐次代入}{求 $(1+x+x^{2})^{1/2}$ 展开式中 $x^{3}$ 项的系数。}
\qi{M4.13}{压轴 · 幂指·二阶}{求 $\lim\limits_{x\to0}\left(\dfrac{(1+x)^{1/x}}{\ee}\right)^{1/x}$。}
\qi{M4.14}{压轴 · 幂指·泰勒}{求 $\lim\limits_{x\to0}\left(\dfrac{\sin x}{x}\right)^{1/x^{2}}$。}
\section{极限计算与间断点（8 题）}
\noindent{\small\color{blue!35!black} B 级 · 加减禁等价替换 · 间断 10/17 套}\\[0.3em]
\noindent\textbf{方法卡}\ 加减式先合并/有理化/展开到第一个非零阶；间断点先分左右，再判可去/跳跃/第二类。\\[0.5em]
\qi{M5.1}{基础 · 加减禁替换}{求 $\lim\limits_{x\to0}\dfrac{x-\sin x}{x^{3}}$。}
\qi{M5.2}{进阶 · 数列·幂指}{求 $\lim\limits_{n\to\infty}\dfrac{n(n^{1/n}-1)}{\ln n}$。}
\qi{M5.3}{进阶 · 数列·二阶}{求 $\lim\limits_{n\to\infty}(1+\tfrac1n)^{n^{2}}\ee^{-n}$。}
\qi{M5.4}{进阶 · 洛必达失效}{求 $\lim\limits_{x\to+\infty}\dfrac{x-\sin x}{x+\sin x}$。}
\qi{M5.5}{进阶 · 洛必达失效·对照}{求 $\lim\limits_{x\to0}\dfrac{x-\sin x}{x+\sin x}$。}
\qi{M5.6}{进阶 · 间断·绝对值}{求 $f(x)=\dfrac{x^{2}-1}{|x-1|}$ 在 $x=1$ 的左右极限并判型。}
\qi{M5.7}{压轴 · 间断·指数}{求 $f(x)=\dfrac1{1+2^{1/x}}$ 在 $x\to0^{-}$、$0^{+}$ 的极限，并判间断类型。}
\qi{M5.8}{压轴 · 幂指·渐近}{求 $\lim\limits_{x\to+\infty}x\left[\left(1+\dfrac1x\right)^{x}-\ee\right]$。}
\section{新压轴：函数方程、迭代与极坐标（8 题）}
\noindent{\small\color{blue!35!black} B 级 · 2024-25 压轴转向 · 训练量最薄}\\[0.3em]
\noindent\textbf{方法卡}\ 函数方程先判条件够不够（不够就给反例）；迭代优先取倒数/取对数化等差；极坐标一律转参数式。\\[0.5em]
\qi{M6.1}{进阶 · 迭代·取倒数}{$a_1=1,\ a_{n+1}=\dfrac{a_n}{1+3a_n}$，求 $\lim n a_n$。}
\qi{M6.2}{基础 · 函数方程}{$f(x+y)=f(x)+f(y)$，$f$ 连续且 $f(1)=3$，求 $f(\sqrt2)$。}
\qi{M6.3}{进阶 · 函数方程·对数型}{$f(xy)=f(x)+f(y)\ (x,y>0)$，$f$ 连续且 $f(\ee)=1$，求 $f\prime(1)$。}
\qi{M6.4}{压轴 · 条件不足·反例}{$f$ 在 $[0,1]$ 连续、$f(0)=f(1)$，是否必存在 $\xi\in(0,1)$ 使 $f\prime(\xi)=0$？给出反例。}
\qi{M6.5}{进阶 · 极坐标切线}{极坐标 $r=1+\sin\theta$，求 $\theta=0$ 处切线斜率。}
\qi{M6.6}{进阶 · 极坐标切线}{极坐标 $r=\ee^{2\theta}$，求 $\theta=0$ 处切线斜率。}
\qi{M6.7}{进阶 · 压缩迭代}{$a_1=1,\ a_{n+1}=\tfrac12(a_n+\tfrac2{a_n})$，求 $\lim a_n$。}
\qi{M6.8}{压轴 · 迭代·渐近阶}{$a_1=1,\ a_{n+1}=\sin a_n$，求 $\lim n a_n^{2}$。}

\section{答案速查（84 行，只给结果）}
{\footnotesize
\setlength{\tabcolsep}{2.5pt}
\renewcommand{\arraystretch}{1.0}
\begin{longtable}{@{}p{1.35cm}p{6.3cm}p{7.6cm}@{}}
\toprule
\textbf{题号} & \textbf{答案} & \textbf{卡点（一句话）} \\
\midrule
\endfirsthead
\toprule
\textbf{题号} & \textbf{答案} & \textbf{卡点（一句话）} \\
\midrule
\endhead
\bottomrule
\endlastfoot
D1 & $y=x+1$ & 先做多项式除法，别把 $b$ 漏掉 \\
D2 & $y=x$ & 分子次数恰比分母高 1 \\
D3 & $\delta=M^{-1/2}$ & 要的是 $>M$，不是 $<\varepsilon$；放缩方向别反 \\
D4 & $\delta=\eps/3$ & 量词行先写全，再化差 \\
D5 & $\dfrac34$；前提 $x\prime_t=2t\ne0$ & 参数方程二阶导要再除以 $x\prime_t$ \\
D6 & $-\dfrac45$；前提 $F_y=x+2y\ne0$ & 先查 $F_y
e0$ 再除 \\
D7 & $\dfrac{21}{64}$ & 通项系数 $\alpha(\alpha-1)(\alpha-2)$，不要乘 $3!$ \\
D8 & $\dfrac13$ & 分母三阶，$\ln(1+x)$ 必须展到三阶 \\
D9 & $\dfrac12$ & 加减式不能逐项等价替换，展到第一个非零阶 \\
D10 & $-1$ 与 $1$；跳跃间断 & $\ee^{1/x}$ 两侧分别趋于 $0$ 与 $+\infty$ \\
D11 & $\dfrac12$ & 取倒数化等差，别直接套压缩映射 \\
D12 & $1$ & 一律转参数式再求 dy/dx，别直接用 r 的导数 \\
M1.1 & $y=x+1$ & 多项式除法得 $b$ \\
M1.2 & $y=2x+1$ & 先约分再判渐近线（$x=-1$ 是可去间断） \\
M1.3 & $y=x$；$x=\pm2$ & 垂直渐近线看分母零点 \\
M1.4 & $y=x$ & 分子分母同次差 1 \\
M1.5 & $y=x+1$ & $\ee^{1/x}=1+\frac1x+o(\frac1x)$ \\
M1.6 & $y=x$ 与 $y=-x$ & $\sqrt{x^2}=|x|$，两端符号不同 \\
M1.7 & $y=x\pm\dfrac\pi2$ & 两端分别取 $\pm\pi/2$ \\
M1.8 & $x=0,\ 2$ & $y\prime=3x(x-2)$ \\
M1.9 & $x=0,\ 2$ & $y\prime=x(2-x)\ee^{-x}$ \\
M1.10 & $x=1$ & $y\prime\prime=6(x-1)$ \\
M1.11 & $x=0,\ 1$ & 要验 $y\prime\prime$ 在两侧变号 \\
M1.12 & $-\dfrac1{2\ee}$ & 驻点 $x=\ee^{-1/2}$，代回时别丢 $\ln x$ \\
M1.13 & $3$ 条：$x=-1,1,2$ & 分母每个单根都算一条 \\
M1.14 & $y=2$ & 同次看首项系数比 \\
M2.1 & $\delta=\eps/3$ & $|f-A|=3|x-1|$ \\
M2.2 & $\delta=\min\{1,\eps/5\}$ & 先限 $|x-2|<1$ 得 $|x+2|<5$ \\
M2.3 & $\delta=\min\{1,\eps/7\}$ & $|x+3|<7$ \\
M2.4 & $\delta=\min\{\tfrac12,\eps/2\}$ & 先限 $|x-1|<1/2$ 得 $|x|>1/2$；分母要上界 \\
M2.5 & $\delta=\eps$ & $|\sin(1/x)|\le1$，别去求 $\sin$ 的极限 \\
M2.6 & $\delta=\min\{1,\sqrt{6\eps}\}$ & 别用洛必达凑 $\delta$，用二阶余项 \\
M2.7 & $X=1/\eps$ & $\eps$-$X$ 的量词行别写成 $\delta$ \\
M2.8 & $X=1/\eps$ & $|f-1|=\frac1{x+1}$ \\
M2.9 & $X=\max\{2,\ 1+3/\eps\}$ & 先保证 $x>1$，$|f-2|=\frac3{|x-1|}$ \\
M2.10 & $\delta=M^{-1/2}$ & 要 $>M$：分母给下界、分子给上界 \\
M2.11 & $\delta=M^{-1/3}$ & 绝对值别丢，双侧都要成立 \\
M2.12 & $\delta=1/M$ & 趋近点是 $1$ 不是 $0$ \\
M2.13 & $\delta=\ee^{-M/2}$ & $2\ln\frac1{|x|}>M$ 反解，别写成 $1/M$ \\
M2.14 & $x_n=\frac1{2n\pi}\Rightarrow1$，$y_n=\frac1{(2n+1)\pi}\Rightarrow-1$ & 两列极限不同即不存在；别只写「振荡」 \\
M3.1 & $\dfrac34$ & 二阶要再除一次 $x\prime_t$ \\
M3.2 & $1$ & $\frac{y\prime_t}{x\prime_t}$，别把 $t$ 当 $x$ \\
M3.3 & $4$ & $x\prime_t=\frac{t}{1+t}$，先化简再代值 \\
M3.4 & $-\dfrac45$ & 隐函数先查 $F_y=x+2y\ne0$ \\
M3.5 & $-1$ & 该点确实在曲线上，先验点再求导 \\
M3.6 & $-\dfrac1\ee$ & $F_y=\ee^{y}+x$，别漏 $x$ \\
M3.7 & $0$ & $y\prime\prime=-\frac{f\prime\prime(y)}{(f\prime(y))^{3}}$ \\
M3.8 & $0$ & 对数求导：$\ln y=\sin x\ln(1+x)$ \\
M3.9 & $a=2,\ b=1$ & 先连续（$b=1$）再可导（$a=2$） \\
M3.10 & $d=2,\ c=3$ & 导数相等给 $d=2$，函数值相等给 $c=3$ \\
M3.11 & $20$ & Leibniz：只有 $k=2$ 项存活，$\binom52\cdot2!=20$ \\
M3.12 & $0$ & $y\prime_t=2\cos2t$ 在该点为 $0$ \\
M3.13 & $\dd y=\dfrac{\dd x}{x}$；前提 $x\ne0$ & 绝对值别丢，前提先写 \\
M3.14 & $y\prime=\dfrac{2x}{\sqrt{1-x^{4}}}$；前提 $|x|<1$ & 链式法则 + $\arcsin$ 定义域 \\
M4.1 & $\dfrac{21}{64}$ & 通项系数，不乘 $3!$ \\
M4.2 & $\dfrac{105}{16}$ & $\alpha=-\frac12$ 的四连乘 \\
M4.3 & $120$ & $x^{6}$ 项系数 $\frac1{3!}$，乘 $6!$ \\
M4.4 & $-120$ & $-\frac{x^{6}}{3!}$ 别写成 $+$ \\
M4.5 & $120$ & $\frac{x^2}{1-x}=x^{2}+x^{3}+\cdots$ \\
M4.6 & $\dfrac1{24}$ & 分母四阶，$\cos x$ 展到四阶 \\
M4.7 & $\dfrac13$ & $\tan x=x+\frac{x^{3}}3+o(x^{3})$ \\
M4.8 & $\dfrac16$ & $\arcsin x=x+\frac{x^{3}}6$ \\
M4.9 & $\dfrac16$ & $\ee^x$ 展到三阶 \\
M4.10 & $-\dfrac14$ & 四阶项系数 $-\frac14$，符号别反 \\
M4.11 & $-60480$ & $x^{9}$ 项 $-\frac1{3!}$，乘 $9!$ \\
M4.12 & $-\dfrac{3}{16}$ & 逐次代入：三项相加得 -3/16，别漏 u 的三次项 \\
M4.13 & $\ee^{-1/2}$ & 取对数：$\frac{\ln(1+x)}x-1\sim-\frac x2$ \\
M4.14 & $\ee^{-1/6}$ & $\ln\frac{\sin x}x\sim-\frac{x^{2}}6$ \\
M5.1 & $\dfrac16$ & 加减式必须展开，不能逐项替换 \\
M5.2 & $1$ & $n^{1/n}-1\sim\frac{\ln n}n$ \\
M5.3 & $\ee^{-1/2}$ & 取对数展开到 $\frac1n$ 项 \\
M5.4 & $1$ & 分子分母同除 $x$，$\frac{\sin x}x\to0$ \\
M5.5 & $0$ & 同一个式子、趋近点不同，答案不同 \\
M5.6 & $2$ 与 $-2$；跳跃间断 & 绝对值分左右拆，别直接约分 \\
M5.7 & $1$ 与 $0$；跳跃间断 & $2^{1/x}$ 两侧分别 $\to0,+\infty$ \\
M5.8 & $-\dfrac\ee2$ & 取对数展开到 $\frac1x$，符号为负 \\
M6.1 & $\dfrac13$ & 取倒数化等差：$1/a_n=3n-2$ \\
M6.2 & $3\sqrt2$ & 连续 + Cauchy $\Rightarrow f(x)=cx$ \\
M6.3 & $1$ & $f(x)=\ln x$，先定函数再求导 \\
M6.4 & 否；反例 $f(x)=|x-\tfrac12|$ & Rolle 需要开区间可导，条件不足要给反例 \\
M6.5 & $1$ & 代参数式求导，别直接用 $r\prime$ \\
M6.6 & $\dfrac12$ & 参数式在 0 点：x 方向导数为 2、y 方向导数为 1 \\
M6.7 & $\sqrt2$ & 单调有界 + 解 $L=\frac12(L+\frac2L)$ \\
M6.8 & $3$ & $\frac1{a_{n+1}^{2}}-\frac1{a_n^{2}}\to\frac13$ \\
\end{longtable}
}
\section{按错因回炉索引}
\begin{itemize}
\item \textbf{渐近线漏 $b$ / 两端符号} $\to$ M1.1--M1.7、D1--D2；回精编本第 7 章七步卡。
\item \textbf{放缩方向反 / 量词写错} $\to$ M2.1--M2.13、D3--D4；回《极限定义证明》三模板与 11/8 执行卡 §6（$\delta=\min\{1/2,1/(2G)\}$、$X=\ee^{M+1}$）。
\item \textbf{漏写前提 / 二阶导少除一次} $\to$ M3.1--M3.14、D5--D6；回冲刺计划 §5-C「每题扫 20 秒」。
\item \textbf{展错阶 / 系数符号} $\to$ M4.1--M4.14、D7--D8；回精编本第 5 章与母函数表。
\item \textbf{加减式误用等价无穷小 / 左右极限} $\to$ M5.1--M5.8、D9--D10；回《函数极限与间断点》§4--§5。
\item \textbf{函数方程条件不足 / 迭代不会化等差} $\to$ M6.1--M6.8、D11--D12；回合订本 p206--p212 详解。
\end{itemize}
\section{最近三天排期（题量加满，不压缩掐表）}
\begin{center}\small
\begin{tabular}{@{}p{2.1cm}p{10.3cm}p{2.9cm}@{}}
\toprule
\textbf{日期} & \textbf{主线} & \textbf{题量} \\
\midrule
11/9 下午 & 下午执行卡不变：模拟卷 A 120 min 掐表 + 50 min 订正 + G1/G2/G4/G9/G5 & 78 \\
11/9 晚间 & \textbf{19:30--19:50 诊断卷 12 题}（20 min 掐表 + 6 行记录）；19:50 起 G6--G8、G10--G11、W1--W5 共 94 题；22:40 星号题复盘 & 12 + 94 \\
11/10 & 早段 G12（C 卷 18 题按单题）+ 模拟卷 B 120 min 掐表 + 订正；下午/晚间按诊断实测顺序做 \textbf{定向加练 72 题}（A 级模块优先）；晚上错题总复盘 + 考场时间分配表 & 36 + 72 + 回炉 \\
11/11 考前 & 只看 5 张手抄卡 + 证明模板卡 + 三种量词行 + \textbf{诊断卷错题重做}；不做新题、不掐表 & 0 新题 \\
\bottomrule
\end{tabular}
\end{center}
\noindent{\footnotesize 今日合计约 184 题；库内题量 \textbf{217 $\to$ 301 题}。A 卷 120 min 掐表、50 min 订正、组末反馈、23:00 停笔 / 23:30 睡一律不压缩；61 页与 74 页阅读包绝对页码未改。答案与题面由 \texttt{build\_weakness\_drill.py} 统一生成，复核全绿后才出 PDF。}
\end{document}
